\frontmatter
\titlepage[crests/officialmono]
\tableofcontents
\clearforchapter\listoffigures
\clearforchapter\listoftables

\frontchapter{Acknowledgements}
This draft reserves space to acknowledge the project supervisor, participants in the planned user study, and anyone who reviewed the software or dissertation. The final wording should be completed by the author before submission.

\frontchapter{Declaration}
This thesis is submitted to the University of Warwick in support of my application for the degree of \thequalification. It has been composed by myself and has not been submitted in any previous application for any degree. The software, test data and analysis presented in this dissertation were produced as part of the project unless a source is explicitly cited. No participant results are reported: the human-comparison protocol is prepared, but recruitment and data collection remain future work. This statement must be reviewed and personalised by the author before submission, including any required declaration concerning the use of generative artificial intelligence.

\frontchapter{Abstract}
Visual behaviour trees are widely used to author game-agent decisions, but conventional node--link editors become difficult to inspect as trees grow and as runtime state changes. This dissertation presents the design and evaluation of a Godot 4.6 editor plugin that combines a complete executable behaviour-tree system with independently switchable techniques for large-tree readability and live debugging. The runtime supports ordered, reactive, random, parallel, repeating and timed control flow; actor method invocation; typed blackboard data; decorators; and editor-only execution inspection. The editor implements compact cards, semantic zoom, fisheye focus+context, subtree collapse and focus, search highlighting, overview+detail, stable layout, alternative edge routing, active-path emphasis and failure annotation.

Evaluation used deterministic 31-, 121- and 364-node trees, a playable 241-node enemy tree, automated semantic and GUI tests, and real-\ac{GPU} visual regression. On the playable tree, each of six display conditions was measured 30 times after three warm-ups, producing 180 raw observations. Compact cards reduced card area by 57.93\% without hiding a card; semantic overview reduced it by 89.48\%; search dimmed 201 of 202 cards; and focus and collapse reduced the canvas to 32 and 43 cards respectively, with zero overlap in every condition. A fixed 240-step drag comparison reduced median processing time by 26.60\%. These are objective presentation and interaction results, not proof of human comprehension. A 12-participant, 216-trial developer study has therefore been implemented with balanced targets for action location, active-path tracing and decorator editing; no participant outcome is reported before recruitment and data collection.

\frontchapter{Abbreviations}
\begin{description}
    \item[AI] Artificial intelligence
    \item[BT] Behaviour tree
    \item[DOI] Degree of interest
    \item[GPU] Graphics processing unit
    \item[GUI] Graphical user interface
    \item[NPC] Non-player character
    \item[RQ] Research question
    \item[UE] Unreal Engine
\end{description}
